\documentclass[10pt,a4paper]{article}
\usepackage[T1]{fontenc}
\usepackage[utf8]{inputenc}
\usepackage{textcomp}\usepackage{lmodern,amsmath,amssymb,array,tabularx,booktabs,xcolor,enumitem}
\usepackage[margin=20mm]{geometry}
\usepackage[unicode,colorlinks=true,urlcolor=black,linkcolor=black]{hyperref}
\usepackage{xcolor}
\setlength{\parindent}{0pt}
\setlength{\parskip}{5pt}
\setlength{\emergencystretch}{2em}
\setlength{\tabcolsep}{4pt}
\renewcommand{\arraystretch}{1.10}
\newcommand{\hd}[1]{\par\vspace{8pt}\textbf{#1}\par\vspace{3pt}}
\newcommand{\circnum}[1]{\textcircled{\scriptsize #1}}
\newcommand{\note}{\textreferencemark\ }
\newcommand{\red}[1]{\textcolor{red}{#1}}
\newcommand{\under}[2]{\underbrace{#1}_{\text{#2}}}
\newcommand{\paper}[1]{\href{#1}{paper}}
\newcolumntype{Y}{>{\raggedright\arraybackslash}X}
\begin{document}

%%%%%%%%%%%%%%%%%%%%%%%%%%%%%%%%%%%%%%%%%%%%%%%%%
\hd{Problem of Interest}
\textbf{Which cities among the potential candidate cities selected by the company should be allocated robotaxis, and how many?}\\

To organize my thoughts, I tried to set a formula starting from simple models to developing into complex models. \textbf{The main model I suggest is the last model.}

\circnum{1} Basic newsvendor model,\\
\circnum{2} Model with added costs\\
\circnum{3} Model that maximizes profit\\
\circnum{4} multi-penalties.

All of the formulations below determine city selection and vehicle numbers simultaneously.

All demand, throughput, and costs use the same evaluation period. After the evaluation period is set, and the model determines the city/vehicle deployment to apply during that period. \textcolor{red}{(Need to think about how to set the period.)}

%%%%%%%%%%%%%%%%%%%%%%%%%%%%%%%%%%%%%%%%%%%%%%%%%%%%
\hd{1. Notation used in common across every model}

{\small
\begin{tabularx}{\linewidth}{@{}>{\raggedright\arraybackslash}p{0.1\linewidth}>{\raggedright\arraybackslash}p{0.36\linewidth}Y@{}}
\toprule
Symbol & Meaning & How is the value determined? \\ \midrule
$C$ & Set of potential candidate cities & The company's candidate list. \\
$S$ & Set of demand scenarios & Need to be generated \\
$y_c$ & 1 if city $c$ is entered, else 0 & \textbf{Determined by the optimization} \\
$q_c$ & Number of vehicles deployed in city $c$ & \textbf{Determined by the optimization} \\
$D_c^s$ & Potential robotaxi demand in scenario $s$ & \textbf{Estimated from data + assumptions for missing parts} \\
$\kappa_c$ & Number of trips one vehicle can serve during the period & Operational data or assumptions or fixed. \\
$Q$ & Total number of vehicles available for deployment & Company input. If unavailable, set an experimental value \\
$\bar q_c$ & Upper bound on vehicle deployment in each city & Facility and operating limits. $Q$ if there is no separate limit \\
$a_c$ & 1 if operation in the city is feasible at the planning stage & Assessed using permit, ODD, and facility data. Separate assumptions for unverified conditions \\
$B$ & Upper bound on the number of cities that can be entered simultaneously & Company input or an experimental value \textbf{only when an actual restriction exists} \\
$R_c^s$ & Actual throughput & A value calculated based on the deployment outcome in the model \\
$O_c^s$ & Remaining service capacity & A value calculated based on the deployment outcome in the model \\
$U_c^s$ & Unserved potential demand & A value calculated based on the deployment outcome in the model \\
\bottomrule
\end{tabularx}
}

$D_c^s$ is a sampled value of uncertain demand. \textbf{After the scenarios are generated, it is a fixed input during optimization}, and the model does not choose demand.

\hd{2. Common constraints}
\[
\boxed{\begin{aligned}
y_c\le q_c\le\bar q_c y_c &\qquad \forall c\in C,\\
\sum_{c\in C}q_c\le Q,\quad \sum_{c\in C}y_c\le B&\\
y_c\le a_c&\qquad\forall c\in C,\\
y_c\in\{0,1\},\quad q_c\in\mathbb Z_{\ge0}.&
\end{aligned}}
\]
\newpage

{\small
\begin{tabularx}{\linewidth}{@{}>{\raggedright\arraybackslash}p{0.35\linewidth}Y@{}}
\toprule
Constraint & Meaning \\ \midrule
$y_c\le q_c\le\bar q_c y_c$ & Zero vehicles if the city is not entered; at least one vehicle is deployed if it is entered \\
$\sum_c q_c\le Q$ & Total vehicle limit \\
$\sum_{c\in C}y_c\le B$ & Limit on the number of cities selected \\
$y_c\le a_c$ & Enter only cities where operation is feasible \\
\bottomrule
\end{tabularx}
}

\hd{3. Values calculated based on the deployment outcome}
Once demand $D_c^s$ and service capacity $\kappa_c q_c$ are determined, the following three values are calculated based on the deployment outcome; \textbf{they are not variables that need to be investigated separately.}
\[
\begin{aligned}
R_c^s&=\min\{D_c^s,\kappa_c q_c\}&&\text{Actual throughput},\\
O_c^s&=(\kappa_c q_c-D_c^s)^+&&\text{Remaining service capacity},\\
U_c^s&=(D_c^s-\kappa_c q_c)^+&&\text{Unserved potential demand}.
\end{aligned}
\]
Here, $(x)^+=\max\{x,0\}$. All three values are measured in trips.

\hd{\circnum{1} Basic form: Minimize the problems of surplus vehicles and shortages}
\[
\boxed{\min_{y,q}\frac{1}{|S|}\sum_{s\in S}\sum_{c\in C}
\left[\under{O_c^s}{Remaining capacity}+\under{\rho U_c^s}{Loss from shortages}\right]}
\]

{\small
\begin{tabularx}{\linewidth}{@{}>{\raggedright\arraybackslash}p{0.14\linewidth}>{\raggedright\arraybackslash}p{0.4\linewidth}Y@{}}
\toprule
Additional symbol & Meaning & How to set it \\ \midrule
$\rho>0$ & How much more important one unit of shortage is considered than one unit of remaining capacity & \textbf{Researcher assumption or company preference.} E.g., sensitivity experiments comparing 1, 2, and 5 \\
\bottomrule
\end{tabularx}
}

\red{This model has no need of precise company cost data, but excludes entry costs of cities.}\\
\

\hd{\circnum{2} Cost form: Minimize total costs, including entry and vehicle costs}
\[
\boxed{\begin{aligned}
\min_{y,q}\quad&\under{\sum_c F_cy_c}{City entry costs}+\under{\sum_c K_cq_c}{Vehicle deployment and holding costs}\\
&+\frac{1}{|S|}\sum_s\sum_c\left[\under{c_c^+O_c^s}{Loss from remaining capacity}+\under{c_c^-U_c^s}{Opportunity loss from unserved demand}\right]
\end{aligned}}
\]
\newpage

{\small
\begin{tabularx}{\linewidth}{@{}>{\raggedright\arraybackslash}p{0.12\linewidth}>{\raggedright\arraybackslash}p{0.37\linewidth}Y@{}}
\toprule
Additional symbol & Meaning and unit & How to set it / Required data \\ \midrule
$F_c$ & Fixed cost of entering the city, in KRW or dollars & \textbf{Additional data required:} Costs of establishing a base and initial preparation. Explicit assumptions if unavailable \\
$K_c$ & Deployment and holding cost per vehicle over the evaluation period & \textbf{Additional data required:} Vehicle, insurance, fixed operating costs, etc. \\
$c_c^+$ & Additional loss per trip of remaining service capacity & Company data or assumptions. \textbf{Do not double-count the same costs as $K_c$} \\
$c_c^-$ & Loss per trip of missed potential demand & Data or assumptions on contribution margin, customer churn, etc. \\
\bottomrule
\end{tabularx}
}

\hd{\circnum{3} Profit form: Maximize the profit from served rides}
\[
\boxed{\begin{aligned}
\max_{y,q}\quad&\frac{1}{|S|}\sum_s\sum_c\under{m_c R_c^s}{Contribution margin from served rides}\\
&-\under{\sum_cF_cy_c}{Entry costs}-\under{\sum_cK_cq_c}{Fixed vehicle costs}
\end{aligned}}
\]

{\small
\begin{tabularx}{\linewidth}{@{}>{\raggedright\arraybackslash}p{0.12\linewidth}>{\raggedright\arraybackslash}p{0.42\linewidth}Y@{}}
\toprule
Additional symbol & Meaning & How to set it / Required data \\ \midrule
$m_c$ & Revenue from one trip minus the variable costs incurred for that trip & \textbf{Additional data required:} Fares and variable costs such as operation and energy. Assumptions if unavailable \\
\bottomrule
\end{tabularx}
}

If customers cannot be served, the profit that would be earned decreases. \textbf{Adding the same revenue loss again as $c_c^-U_c^s$ may result in double counting.}

This most directly reflects the company's perspective, but \red{fare and variable-cost data are required.} If it is not profitable, the result may be to enter no city.
\end{document}
